\chapter{Scénarios d'utilisation courants}
\minitoc

Ce chapitre déroule les six situations que rencontre un exploitant. Chacune se
lit indépendamment des autres.

\section{Scénario 1 --- Une menace connue est traitée seule}

\textbf{Situation.} Une sonde signale une exfiltration de données par tunnel
DNS depuis un serveur de base de données.

\textbf{Ce que fait la plateforme, sans vous.}

\begin{enumerate}[leftmargin=*]
  \item Elle normalise l'observation et vérifie qu'elle ne l'a pas déjà traitée.
  \item Elle la rapproche du catalogue : \emph{A5, exfiltration de données},
        gravité critique, dangerosité 9,4 sur 10, priorité critique.
  \item Elle cherche les cas comparables dans sa documentation : quatre fiches
        concordent.
  \item Elle applique la conduite à tenir documentée : coupure de la connexion
        sortante, blocage de l'adresse au pare-feu, bascule de la machine en
        réseau de quarantaine.
  \item Elle vous informe.
\end{enumerate}

\textbf{Ce que vous faites.} Vous ouvrez \textbf{Vue d'ensemble} et vous
constatez. Si un geste vous paraît inapproprié, vous l'annulez depuis le détail
de l'intervention : chaque geste appliqué porte un jeton d'annulation qui
permet de le défaire précisément.

\begin{tcolorbox}[colback=green!4!white, colframe=darkgreen,
  title=\textbf{Bon à savoir}, fonttitle=\bfseries, boxrule=0.6pt]
Le délai entre la réception de l'observation et le dernier geste appliqué se
compte en dizaines de millisecondes. Quand vous ouvrez la console, tout est
déjà fait.
\end{tcolorbox}

\section{Scénario 2 --- Une menace inédite se présente}

\textbf{Situation.} Une sonde signale une séquence protocolaire non conforme
sur un automate industriel, suivie d'un transfert sortant massif. Aucun type du
catalogue ne correspond.

\textbf{Ce que fait la plateforme.} Elle ne devine pas. N'ayant aucun cas
comparable documenté, elle s'interdit d'appliquer une conduite à tenir. Mais
elle ne reste pas inerte pour autant : elle bascule sur un \textbf{confinement
de repli} fondé sur les seuls indicateurs observés.

\begin{quote}
\small
« Adresse source externe observée : 185.220.101.44~; destination sortante
externe observée : 91.199.7.12~; compte impliqué : operateur.scada~; port
concerné : 502~; actif de criticité 5/5 exposé : srv-scada-07. »
\end{quote}

De ces observations elle déduit quatre gestes \textbf{entièrement réversibles},
qu'elle applique : blocage de l'adresse, exigence d'authentification renforcée
sur le compte, fermeture du port, sauvegarde immédiate.

Quatre autres gestes seraient pertinents mais ont un \emph{effet durable} :
coupure de la connexion sortante, fermeture des sessions, verrouillage du
compte, mise en quarantaine réseau. Elle ne les exécute pas. Elle les
\textbf{inscrit en attente de votre décision}.

Enfin, elle ouvre une \textbf{fiche de qualification} : une proposition de nom
pour cette menace qu'elle ne sait pas désigner.

\textbf{Ce que vous faites.} Vous traitez les décisions requises
(scénario~3) puis la fiche de qualification (scénario~5).

\section{Scénario 3 --- Trancher sur un geste à effet durable}

\textbf{Situation.} La \textbf{Vue d'ensemble} affiche un bloc « décisions
requises » avec, pour chaque geste : ce qu'il ferait, sur quoi, pourquoi la
plateforme le propose, et ce qui subsisterait après annulation.

\begin{tcolorbox}[colback=green!4!white, colframe=darkgreen,
  title=\textbf{Bon à savoir --- ce n'est pas une validation préalable},
  fonttitle=\bfseries, boxrule=0.6pt]
Ces gestes n'ont jamais été planifiés pour exécution autonome. Le planificateur
les a écartés au moment de la décision, parce que le catalogue les déclare à
effet durable. Rien n'attend ici pour partir : ce qui attend, c'est une
décision \emph{humaine} sur un geste que la plateforme a refusé de s'autoriser.
\end{tcolorbox}

Saisissez le motif de votre décision dans le champ prévu — il est consigné au
journal et sans lui rien ne part — puis choisissez.

\subsection{Confirmer}

Vous assumez l'effet durable. La plateforme exécute le geste par la voie
nominale : contrôle avant vol, journal, jeton d'annulation, surveillance de
l'effet. Une action confirmée par un humain n'est pas de seconde classe :
elle reste annulable comme les autres.

\subsection{Je m'en charge}

Vous intervenez vous-même sur l'équipement. La plateforme n'exécute rien et
pose un statut de \textbf{prise en charge}. Le dossier \textbf{ne se referme
pas} : il reste visible à votre nom jusqu'à ce que vous rendiez compte de ce
qui a été fait, par le bouton \textbf{Rendre compte et clore}.

\begin{tcolorbox}[colback=green!4!white, colframe=darkgreen,
  title=\textbf{Bon à savoir}, fonttitle=\bfseries, boxrule=0.6pt]
Ce fonctionnement en deux temps est délibéré. Un engagement que plus personne
ne voit est exactement le défaut que l'alerte persistante corrige.
\end{tcolorbox}

\subsection{Écarter}

Vous jugez le geste inutile ou disproportionné. \textbf{Il ne sera jamais
exécuté}, ni sous ce nom ni sous un autre.

Mais votre refus dit que \emph{ce geste-là} ne convient pas — il ne rend pas la
menace inoffensive. La plateforme cherche donc ce qu'elle sait encore faire,
et vous répond. Deux cas, arbitrés par la dangerosité de l'intervention :

\begin{table}[h]
\centering
\begin{tabular}{@{}p{3.4cm}p{10.4cm}@{}}
\toprule
\textbf{Dangerosité} & \textbf{Ce que fait la plateforme} \\
\midrule
Sous 7 sur 10 &
\textbf{Surveillance rapprochée.} Aucun geste n'est posé : votre refus
s'applique tel quel. L'actif reste sous l'œil, l'intervention demeure ouverte,
et un geste plus léger vous est proposé — appliqué seulement si vous le
demandez. \\[8pt]
7 sur 10 et au-delà &
\textbf{Confinement de substitution.} La plateforme applique un geste plus
léger servant le même objectif, entièrement réversible, que vous pouvez annuler
à tout moment. \\
\bottomrule
\end{tabular}
\caption{L'arbitrage après un refus}
\end{table}

\textbf{Exemple concret.} Vous écartez la mise en quarantaine réseau d'un
automate, au motif qu'il est en production. La plateforme vous propose la
\emph{bascule en réseau de quarantaine} : même objectif — couper la machine du
reste du réseau — mais entièrement réversible là où la première ne l'était
qu'à moitié. Le geste que vous avez écarté n'est pas rejoué.

\begin{tcolorbox}[colback=green!4!white, colframe=darkgreen,
  title=\textbf{Bon à savoir --- d'où vient la proposition},
  fonttitle=\bfseries, boxrule=0.6pt]
La plateforme rattache chaque geste de son catalogue à un \emph{objectif de
confinement}, et calcule pour chacun son \emph{engagement} : ce qu'il coûte
s'il se révèle inutile. Une alternative n'est retenue que si elle sert le même
objectif, s'annule entièrement, et engage \textbf{strictement moins} que le
geste refusé. Proposer plus engageant après un refus reviendrait à passer
outre votre décision.
\end{tcolorbox}

\section{Scénario 4 --- Annuler un geste appliqué}

\textbf{Situation.} La plateforme a bloqué une adresse qui se révèle
légitime — un prestataire, un site distant du Ministère.

Depuis le détail de l'intervention, chaque geste exécuté porte un bouton
d'annulation. L'annulation emploie le jeton rendu par l'équipement au moment
de l'application : elle défait \emph{précisément} ce geste, et rien d'autre.

\begin{tcolorbox}[colback=green!4!white, colframe=darkgreen,
  title=\textbf{Bon à savoir}, fonttitle=\bfseries, boxrule=0.6pt]
Vous n'êtes pas seul à pouvoir annuler. Après un délai de surveillance, la
plateforme compare l'état de la machine avant et après son intervention. Si le
confinement a dégradé le service qu'il protégeait, elle l'annule d'elle-même.
Les annulations autonomes et manuelles sont distinguées dans les rapports.
\end{tcolorbox}

\section{Scénario 5 --- Nommer une menace inédite}

\textbf{Situation.} Une fiche de qualification attend dans la
\textbf{Vue d'ensemble}.

La fiche porte un nom proposé, une famille déduite, une gravité, une
dangerosité, et surtout la \textbf{signature} : les indicateurs qui ont été
observés, et qui serviront à reconnaître la prochaine occurrence.

\begin{enumerate}[leftmargin=*]
  \item Lisez le nom proposé. Il décrit ce qui a été constaté, pas un
        diagnostic.
  \item Corrigez-le s'il ne porte pas le sens métier attendu. Dans notre
        exemple, « Communication sortante vers une destination externe non
        référencée » devient « Détournement de protocole industriel Modbus ».
  \item \textbf{Valider et cataloguer.} Le type rejoint le catalogue appris
        sous un code de la forme \texttt{L01}.
\end{enumerate}

À l'occurrence suivante — même forme d'attaque, machine et adresse
différentes — la plateforme reconnaît le type et l'intervention est
correctement qualifiée et priorisée.

\begin{tcolorbox}[colback=red!4!white, colframe=darkred,
  title=\textbf{Limitation connue --- apprendre à nommer, pas à répondre},
  fonttitle=\bfseries, boxrule=0.6pt]
Adopter une fiche enrichit la \emph{classification} : l'incident suivant est
reconnu, sa dangerosité entre dans le score de risque, son rang au portefeuille
est juste. Mais la \emph{réponse} continue de venir du confinement de repli,
faute de conduite à tenir écrite pour ce type appris. Fermer complètement la
boucle demande qu'un analyste rédige un playbook — geste délibérément laissé
humain.
\end{tcolorbox}

\section{Scénario 6 --- La plateforme s'est trompée trop souvent}

\textbf{Situation.} Le bandeau affiche \textbf{Autonomie suspendue}.

Le \emph{coupe-circuit} s'est déclenché : sur une fenêtre de dix minutes, le
nombre d'annulations ou d'échecs a dépassé le seuil configuré. La plateforme
continue de recevoir, de classifier et de journaliser, mais elle
\textbf{n'exécute plus rien}.

\textbf{Ce que vous faites.}

\begin{enumerate}[leftmargin=*]
  \item Ouvrez le \textbf{Journal d'audit} et filtrez sur les annulations pour
        comprendre ce qui a dérivé.
  \item Corrigez la cause : une consigne de réponse trop large, un playbook
        inadapté, un actuateur en panne.
  \item Réarmez le coupe-circuit depuis les \textbf{Réglages} — geste réservé à
        l'administrateur, et inscrit au journal avec son auteur et son motif.
\end{enumerate}

\begin{tcolorbox}[colback=orange!5!white, colframe=orange!70!black,
  title=\textbf{Attention}, fonttitle=\bfseries, boxrule=0.6pt]
Ne réarmez pas sans avoir traité la cause. Le coupe-circuit s'est déclenché
parce que la plateforme se corrigeait en boucle : le réarmer sans rien changer
la ramène dans le même état, et le journal s'emplit d'annulations qui masquent
les vraies.
\end{tcolorbox}
